%----------------------------------------------------------
% 제18장. 공공부문 AI 거버넌스
%----------------------------------------------------------

\chapter{공공부문 AI 거버넌스}
\label{ch:public_governance}

공공부문 AI는 민간 AI와 근본적으로\cite{berryhill2019hello} 다른 거버넌스 원칙을 요구한다. 민간 기업의 AI가 효율성 수익성을 최우선으로 추구한다면, 공공 AI는 공공성\cite{wirtz2019artificial}, 형평성, 민주적 책임성\cite{coglianese2017regulating}이라는 헌법적 가치를 실현해야 한다. 공공 AI의 결정은 시민의 기본권에 직접적 영향을 미치며, 그 영향의 범위와 깊이는 민간 서비스와 비교할 수 없을 정도로 광범위하다.

공공부문 AI 거버넌스의 핵심 과제는 세 가지 긴장 관계를 관리하는 것이다. 첫째, \textbf{행정 효율성 절차적 정당성} 간의 균형이다. AI 자동화는 행정 처리 속도를 획기적으로 높일 수 있으나, 적법 절차와 이의제기 권리를 훼손해서는 안 된다. 둘째, \textbf{데이터 활용과 프라이버시 보호} 간의 균형이다. 공공데이터\cite{janssen2020data}의 AI 활용은 정책 품질을 향상시키지만, 시민 감시 우려를 야기\cite{eubanks2018automating}할 수 있다. 셋째, \textbf{기술적 복잡성 민주적 통제} 간의 균형이다. AI 시스템의 기술적 불투명성\cite{arrieta2020explainable}이 시민과 의회의 제어력을 약화\cite{shneiderman2020bridging}시켜서는 안 된다.

본 장에서는 공공 AI의 특수성을 분석하고, 국내 관련 법제를 검토하며, 공공 AI 조달\cite{kuziemski2020ai}과 검증의 실무 체계를 제시한다. 또한 행정 자동화\cite{brown2019government}, 지능형 교통, 민원 상담 AI 등 주요 응용 영역의 거버넌스 사례를 심층적으로 다룬다.

%----------------------------------------------------------
\section{공공 AI의 특수성}
\label{sec:18.1}

\subsection{공공성과 형평성의 원칙}
\label{sec:18.1.1}

공공 AI는 \textbf{공공성(Publicness)}의 원칙에 기반해야 한다. 공공성은 단순히 '정부가 운영한다'는 의미를 넘어, AI 시스템이 모든 시민에게 보편적으로 접근 가능하고, 특정 집단을 차별하지 않으며, 공익을 위해 작동해야 함을 의미한다.

\begin{table}[htbp]
\centering
\caption{공공 AI의 공공성 원칙}
\label{tab:18_1}
\begin{tabularx}{\textwidth}{p{2cm}Xp{5cm}}
\toprule
\textbf{원칙} & \textbf{내용} & \textbf{거버넌스 구현} \\
\midrule
보편성 & 모든 시민이 차별 없이 서비스 이용 가능 & 디지털 접근성 보장, 오프라인 대안 제공 \\
형평성 & 사회적 약자에 대한 불리한 결과 방지 & 취약계층 영향 평가, 편향\cite{mehrabi2021survey} 모니터링 \\
투명성 & 의사결정 과정 공개와 설명 & 알고리즘\cite{veale2019administration} 공개, 설명 요청권 보장 \\
책임성 & 결과에 대한 명확한 책임 귀속 & 인간 최종 결정, 구제 절차 마련 \\
참여성 & 시민의 의견 수렴과 제어 & 공청회, 시민 감사, 의회 보고 \\
\bottomrule
\end{tabularx}
\end{table}

\textbf{형평성(Equity)}은 공공 AI의 핵심 요건이다. 민간 AI에서 공정성이 '차별 금지'의 소극적 의미라면, 공공 AI에서 형평성은 '사회적 약자 배려'라는 적극적 의미를 포함한다. 예컨대 복지 수급 자격 심사 AI가 동일한 기준을 적용하더라도, 그 결과가 특정 지역이나 계층에 불리하게 작용한다면 형평성 원칙에 위배된다.

\begin{lstlisting}[language=Python, caption={공공 AI Equity Assessment and Vulnerable Group Impact Analysis Framework}, label={lst:equity_eval}]
from dataclasses import dataclass
from enum import Enum
import pandas as pd
import numpy as np

class VulnerableGroup(Enum):
 ELDERLY = " "; DISABLED = " "; LOW_INCOME = " "
 RURAL = " "; DIGITAL_VULNERABLE = " "

class PublicAIEquityEvaluator:
 """ Equity Assessment : Metric """
 def __init__(self, disparity_threshold=0.15, critical_threshold=0.30):
 self.disparity_threshold = disparity_threshold
 self.critical_threshold = critical_threshold

 def evaluate_equity(self, data, group_col, outcome_col, groups):
 results = []
 general_mask = ~data[group_col].isin([g.value for g in groups])
 general_val = data[general_mask][outcome_col].mean()

 for group in groups:
 group_val = data[data[group_col] == group.value][outcome_col].mean()
 disparity = (general_val - group_val) / general_val if general_val != 0 else 0

 status = " "
 if abs(disparity) > self.critical_threshold: status = " "
 elif abs(disparity) > self.disparity_threshold: status = " "

 results.append({"group": group.value, "disparity": disparity, "status": status})
 return results
\end{lstlisting}

\paragraph{공공성 원칙 심화: 해외 실패 사례로부터의 교훈}

공공 AI 거버넌스에서 공공성 원칙이 왜 중요한지를 보여주는 대표적 해외 사례가 존재한다. 이 사례들은 공공 AI가 형평성과 투명성 원칙을 위반했을 때 발생할 수 있는 심각한 결과를 경고한다.

\textbf{네덜란드 SyRI(System Risk Indication) 사건}: 네덜란드 정부는 2014년부터 복지 사기를 탐지하기 위해 SyRI 시스템을 운영했다. 이 시스템은 세금, 사회보험, 고용 등 다양한 정부 데이터베이스를 결합하여 '리스크 프로파일'을 생성했다. 그러나 2020년 헤이그 지방법원은 SyRI가 유럽인권협약 제8조(사생활 보호권)를 위반한다고 판결했다. 핵심 쟁점은 세 가지였다. 첫째, 알고리즘의 작동 방식이 시민에게 전혀 공개되지 않았다. 둘째, 시스템이 저소득 지역에 집중 적용되어 사회적 약자를 차별적으로 감시하는 결과를 초래했다. 셋째, 시민이 자신의 리스크 점수를 확인하거나 이의를 제기할 방법이 없었다.

\textbf{호주 Robodebt 사건}: 호주 정부는 2016년부터 복지 수급자의 소득 신고를 자동 대조하여 과다 수급액을 산정하고 자동으로 채무 통지를 발송하는 시스템을 운영했다. 그러나 연간 소득을 단순 균등 분할하여 주급과 비교하는 알고리즘의 구조적 결함으로 인해, 약 47만 건의 부당한 채무 통지가 발송되었다. 2020년 호주 연방법원은 이를 불법으로 판결하였고, 정부는 약 19억 호주달러의 배상금을 지급해야 했다.

\begin{table}[htbp]
\centering
\caption{공공 AI 해외 실패 사례와 거버넌스 교훈}
\label{tab:18_failures}
\begin{tabularx}{\textwidth}{p{2.5cm}p{2cm}Xp{4cm}}
\toprule
\textbf{사례} & \textbf{국가} & \textbf{핵심 문제} & \textbf{거버넌스 교훈} \\
\midrule
SyRI & 네덜란드 & 알고리즘 불투명, 차별적 적용, 이의제기 불가 & 투명성 레지스터, 영향평가\cite{raji2020closing} 의무화, 시민 접근권 \\
Robodebt & 호주 & 알고리즘 설계 결함, 인간 검토 부재 & 자동화 범위 제한, 인간 개입 보장 \\
COMPAS & 미국 & 인종별 재범 예측 편향 & 공정성 감사, 보호 속성 영향 분석 \\
A-levels & 영국 & 학교별 과거 성적 기반 차별 & 알고리즘 사전 영향 시뮬레이션 \\
\bottomrule
\end{tabularx}
\end{table}

이러한 해외 사례들은 한국의 공공 AI 거버넌스 설계에 중요한 시사점을 제공한다. 특히 (1) 알고리즘의 사전 공개 또는 등록 의무, (2) 취약계층에 대한 차등 영향 평가, (3) 자동화 범위의 법적 제한, (4) 시민의 이의제기권과 구제 절차 확보가 필수적이다.

\subsection{시민 신뢰와 민주적 책임성}
\label{sec:18.1.2}

공공 AI에 대한 시민 신뢰는 단순한 기술 신뢰를 넘어 제도적 신뢰에 기반한다. 시민들은 AI 시스템 자체보다 그 시스템을 운영하는 정부 기관, 감독하는 의회, 구제하는 사법부를 신뢰해야 한다. 따라서 공공 AI 거버넌스는 기술적 안전성뿐 아니라 민주적 제어 메커니즘의 구축을 포함해야 한다.

\textbf{민주적 책임성(Democratic Accountability)}은 세 차원에서 구현된다(표 \ref{tab:18_2}). 또한 공공 AI의 민주적 제어를 위해서는 알고리즘의 적절한 수준의 공개가 필요하다. 완전한 소스 코드 공개는 보안 위험이 있으나, 의사결정 로직의 핵심 요소는 시민이 이해할 수 있는 형태로 공개되어야 한다.

\begin{table}[htbp]
\centering
\caption{공공 AI의 민주적 책임성 구조}
\label{tab:18_2}
\begin{tabularx}{\textwidth}{p{2cm}p{3cm}X}
\toprule
\textbf{차원} & \textbf{주체} & \textbf{메커니즘} \\
\midrule
행정적 책임 & 담당 기관 & 시스템 운영, 품질 관리, 감사 대응 \\
정치적 책임 & 의회/지방의회 & 예산 제어, 정책 감독, 국정감사 \\
법적 책임 & 법원/헌법재판소 & 권리 구제, 위법 결정, 행정소송 \\
\bottomrule
\end{tabularx}
\end{table}

\paragraph{알고리즘 투명성 레지스터(Algorithm Transparency Register)}

알고리즘 투명성 레지스터는 정부가 사용하는 AI/알고리즘 시스템의 목록과 상세 정보를 공개하는 제도이다. 이는 시민이 어떤 알고리즘이 자신에게 영향을 미칠 수 있는지 알 권리를 보장하는 핵심 메커니즘이다.

\textbf{네덜란드 알고리즘 레지스터}: 네덜란드는 SyRI 사건의 교훈을 바탕으로 2022년 정부 차원의 알고리즘 레지스터(Algoritmeregister)를 구축하였다. 이 레지스터에는 각 알고리즘의 목적, 사용 데이터, 작동 방식, 리스크 등급, 영향평가 결과, 인간 개입 수준이 기록되며, 시민 누구나 웹사이트를 통해 열람할 수 있다.

\textbf{영국 알고리즘 투명성 기록 표준(Algorithmic Transparency Recording Standard, ATRS)}: 영국 중앙디지털데이터사무국(CDDO)은 2022년 ATRS를 발표하여 정부 기관이 알고리즘 도구의 정보를 표준화된 형식으로 공개하도록 하였다. 이 표준은 두 단계로 구성된다: Tier 1(기본 정보)은 시스템 이름, 목적, 개요를 포함하고, Tier 2(상세 정보)는 데이터 소스, 리스크 평가, 영향 분석, 인간 개입 수준을 포함한다.

\begin{table}[htbp]
\centering
\caption{알고리즘 투명성 레지스터 필수 기재 항목}
\label{tab:18_register}
\begin{tabularx}{\textwidth}{p{3cm}Xp{4cm}}
\toprule
\textbf{기재 항목} & \textbf{내용} & \textbf{공개 수준} \\
\midrule
시스템 명칭 & AI 시스템의 공식 명칭 및 버전 & 완전 공개 \\
사용 목적 & 어떤 행정 업무에 어떻게 사용되는지 & 완전 공개 \\
사용 데이터 & 입력 데이터의 종류 및 출처 & 완전 공개 \\
작동 방식 & 핵심 알고리즘 로직 설명 & 요약 공개 \\
리스크 등급 & 시민 권리 영향 수준 & 완전 공개 \\
영향평가 결과 & 공정성/형평성 평가 요약 & 요약 공개 \\
인간 개입 수준 & 자동/반자동/보조 구분 & 완전 공개 \\
이의제기 절차 & 시민이 결과에 이의를 제기하는 방법 & 완전 공개 \\
최종 검증일 & 마지막 검증/감사 수행 일자 & 완전 공개 \\
담당 기관/부서 & 운영 책임 조직 및 연락처 & 완전 공개 \\
\bottomrule
\end{tabularx}
\end{table}

\subsection{행정 효율성 인권 보호의 균형}
\label{sec:18.1.3}

공공 AI 도입의 주된 동기는 행정 효율성 향상이다. 그러나 효율성 추구가 시민의 기본권을 침해해서는 안 된다. \textbf{비례성 원칙(Principle of Proportionality)}에 따라 AI 도입으로 인한 기본권 제한은 정당한 목적 달성에 적합하고, 필요 최소한의 수단이어야 하며, 공익적 기대 이익이 침해되는 사익보다 커야 한다.

%----------------------------------------------------------
\section{공공 AI 규제와 정책}
\label{sec:18.2}

\subsection{지능정보화 기본법과 AI 활용}
\label{sec:18.2.1}

'지능정보화 기본법'은 공공부문 AI 활용의 기본 법적 근거를 제공한다 \cite{korea2020intel}. 제56조(안전성 확보)에 따라 공공 AI는 안전성 평가와 지속적 모니터링 의무를 부담하며, 제62조(지능정보사회 윤리)는 인간의 존엄을 보장하는 윤리 정립 의무를 명시한다.

\subsection{행정기본법상 자동화된 처분}
\label{sec:18.2.2}

2021년 시행된 '행정기본법' 제20조는 자동적 처분에 관한 명시적 근거를 마련했다 \cite{korea2021admin}. 핵심은 \textbf{법률 유보 원칙}(반드시 법률 근거 필요), \textbf{기속 행위 제한}(재량이 있는 처분은 완전 자동화 금지), \textbf{절차적 보장}(의견 제출 및 불복 고지 필수)이다.

\paragraph{행정기본법 제20조 심층 분석}

행정기본법 제20조는 공공 AI 거버넌스의 법적 토대를 구성하는 핵심 조항이다. 이 조항의 실무적 함의를 깊이 분석한다.

\textbf{자동적 처분의 유형}: 행정기본법 제20조가 적용되는 자동적 처분은 크게 세 가지 유형으로 구분된다.

\begin{table}[htbp]
\centering
\caption{행정기본법 제20조 자동적 처분 유형과 AI 적용 범위}
\label{tab:18_admin_types}
\begin{tabularx}{\textwidth}{p{2.5cm}Xp{3cm}p{2.5cm}}
\toprule
\textbf{처분 유형} & \textbf{설명} & \textbf{AI 적용 예시} & \textbf{자동화 허용 범위} \\
\midrule
기속 행위 & 법령상 요건 충족 시 반드시 처분 & 연금 수급 자격, 면허 갱신 & 완전 자동화 가능 \\
제한적 재량 & 일정 범위 내 재량 존재 & 과태료 부과, 보조금 산정 & 반자동화(AI 보조) \\
광범위 재량 & 담당자 판단 필수 & 허가/인가 심사, 징계 & 자동화 금지 \\
\bottomrule
\end{tabularx}
\end{table}

\textbf{법적 쟁점}: 행정기본법 제20조의 적용에는 여러 법적 쟁점이 존재한다.

첫째, \textbf{기속 행위와 재량 행위의 경계} 문제이다. 현실에서 순수한 기속 행위는 드물며, 대부분의 행정 처분에는 크든 작든 재량 요소가 포함된다. 복지 수급 자격 심사에서 '소득'의 범위를 어디까지로 볼 것인지, '가구'의 범위를 어떻게 판단할 것인지 등은 실질적으로 재량의 영역이다.

둘째, \textbf{AI의 '판단'과 '처분'의 구분} 문제이다. AI가 정보를 분석하여 결과를 제시하는 것(판단 보조)과 AI가 직접 행정 처분을 내리는 것(자동적 처분)은 법적으로 구분되어야 한다. 전자는 행정기본법 제20조의 적용을 받지 않으나, 후자는 엄격한 요건을 충족해야 한다.

셋째, \textbf{절차적 권리 보장}의 실효성이다. 자동적 처분을 받은 시민에게 의견 제출과 불복 절차가 형식적으로 보장되더라도, AI 시스템의 의사결정 과정을 이해하지 못하면 실질적 권리 행사가 어렵다. 따라서 설명가능성은 절차적 권리의 전제 조건이다.

\begin{lstlisting}[language=Python, caption={Administrative Basic Act Article 20 Compliance Checker for Public AI}, label={lst:admin_compliance}]
from dataclasses import dataclass
from typing import Dict, List, Optional, Tuple
from enum import Enum
from datetime import datetime

class DispositionType(Enum):
    BOUND = "bound_act"        # (fully automatable)
    LIMITED_DISCRETION = "limited_discretion"  # (AI-assisted)
    BROAD_DISCRETION = "broad_discretion"      # (no automation)

class AutomationLevel(Enum):
    FULL_AUTO = "full_automation"
    SEMI_AUTO = "semi_automation"
    AI_ASSISTED = "ai_assisted"
    HUMAN_ONLY = "human_only"

@dataclass
class DispositionAssessment:
    disposition_name: str
    disposition_type: DispositionType
    legal_basis: str
    allowed_automation: AutomationLevel
    requires_explanation: bool
    requires_appeal_notice: bool
    human_review_required: bool
    compliance_issues: List[str]

class Article20ComplianceChecker:
    """Check AI system compliance with Administrative Basic Act Art.20"""

    AUTOMATION_RULES = {
        DispositionType.BOUND: AutomationLevel.FULL_AUTO,
        DispositionType.LIMITED_DISCRETION: AutomationLevel.SEMI_AUTO,
        DispositionType.BROAD_DISCRETION: AutomationLevel.HUMAN_ONLY,
    }

    def __init__(self):
        self.assessments: List[DispositionAssessment] = []

    def assess_disposition(self, name: str, dtype: DispositionType,
                           legal_basis: str,
                           current_automation: AutomationLevel) -> DispositionAssessment:
        """Assess whether current automation level is legally compliant"""
        allowed = self.AUTOMATION_RULES[dtype]
        issues = []

        # Check automation level compliance
        auto_order = [AutomationLevel.HUMAN_ONLY, AutomationLevel.AI_ASSISTED,
                     AutomationLevel.SEMI_AUTO, AutomationLevel.FULL_AUTO]
        if auto_order.index(current_automation) > auto_order.index(allowed):
            issues.append(
                f"Automation level '{current_automation.value}' exceeds "
                f"allowed level '{allowed.value}' for {dtype.value}"
            )

        # Legal basis check
        if not legal_basis or legal_basis.strip() == "":
            issues.append("Missing legal basis for automated disposition")

        # Discretion check
        human_review = dtype != DispositionType.BOUND

        assessment = DispositionAssessment(
            disposition_name=name,
            disposition_type=dtype,
            legal_basis=legal_basis,
            allowed_automation=allowed,
            requires_explanation=True,  # Always required for public AI
            requires_appeal_notice=True,  # Always required
            human_review_required=human_review,
            compliance_issues=issues
        )
        self.assessments.append(assessment)
        return assessment

    def generate_compliance_report(self) -> Dict:
        """Generate Article 20 compliance report"""
        total = len(self.assessments)
        compliant = sum(1 for a in self.assessments if not a.compliance_issues)

        return {
            "report_date": datetime.now().isoformat(),
            "total_dispositions": total,
            "compliant": compliant,
            "non_compliant": total - compliant,
            "issues": [
                {"disposition": a.disposition_name, "issues": a.compliance_issues}
                for a in self.assessments if a.compliance_issues
            ],
            "overall_status": "COMPLIANT" if compliant == total else "ACTION_REQUIRED"
        }
\end{lstlisting}

\subsection{공공데이터 활용과 AI의 교차}
\label{sec:18.2.3}

공공데이터의 AI 활용은 정책 품질을 높이지만, 목적 외 이용과 재식별 위험이라는 과제를 안고 있다. 특히 최근 \textbf{공공 마이데이터(Public MyData)} 정책이 확대되면서, 시민이 직접 자신의 행정 정보를 AI 서비스에 제공하여 맞춤형 혜택을 받는 구조가 확산되고 있다.

공공-민간 데이터 결합 시에는 개인정보보호법상 \textbf{가명처리}와 \textbf{데이터 결합 전문기관}의 역할을 준수해야 하며, 데이터 패브릭(Data Fabric) 기반의 통합 가상화 계층을 통해 물리적 데이터 이동 없이 보안을 유지하며 AI 학습을 수행하는 거버넌스 아키텍처가 권장된다.

\paragraph{공공데이터 AI 활용 상세 프로세스}

공공데이터를 AI에 활용하는 과정은 데이터 식별, 가명처리, 결합, 활용, 폐기의 라이프사이클로 관리되어야 한다.

\begin{table}[htbp]
\centering
\caption{공공데이터 AI 활용 라이프사이클과 거버넌스 요건}
\label{tab:18_data_lifecycle}
\begin{tabularx}{\textwidth}{p{2cm}Xp{4cm}}
\toprule
\textbf{단계} & \textbf{주요 활동} & \textbf{거버넌스 요건} \\
\midrule
데이터 식별 & 필요 데이터 식별, 민감도 분류 & 목적 명확화, 최소 수집 원칙 \\
가명처리 & K-익명성, L-다양성 적용 & 적정성 평가, 재식별 위험 검증 \\
데이터 결합 & 결합 전문기관 경유 & 결합키 관리, 반출 심사 \\
AI 학습 & 모델 학습/검증 & 목적 외 이용 금지, 편향 점검 \\
서비스 운영 & AI 추론/의사결정 & 설명 제공, 이의제기 보장 \\
데이터 폐기 & 목적 달성 후 파기 & 파기 증명서 발행, 감사 추적 \\
\bottomrule
\end{tabularx}
\end{table}

\textbf{데이터 결합의 거버넌스}: 서로 다른 공공기관의 데이터를 결합하여 AI에 활용하는 경우, 개인정보보호법 제28조의3에 따라 \textbf{데이터 결합 전문기관}을 경유해야 한다. 결합 전문기관은 결합키 관리, 가명처리 적정성 검증, 반출 데이터 심사를 수행한다. 공공 AI의 경우 이 과정에서 추가적으로 (1) 결합 목적의 공익성 검증, (2) 결합 데이터의 재식별 리스크 평가, (3) AI 학습 완료 후 결합 데이터 파기 확인이 요구된다.

\textbf{공공 마이데이터(MyData)}: 공공 마이데이터는 시민이 자신의 행정 정보(소득, 재산, 건강, 교육 등)를 통합 조회하고, 본인 동의 하에 제3자에게 제공할 수 있는 체계이다. AI 서비스와 결합되면 맞춤형 복지 안내, 개인화된 행정 서비스 등이 가능하나, 동의의 진정성(Informed Consent), 데이터 최소 수집, 목적 외 이용 금지 등의 거버넌스 원칙이 철저히 준수되어야 한다.

%----------------------------------------------------------
\section{시민 참여 메커니즘}
\label{sec:18.2.4}

\subsection{시민 참여의 유형과 단계}

공공 AI 거버넌스에서 시민 참여는 민주적 정당성을 확보하기 위한 핵심 요소이다. 시민 참여는 단순 정보 공개부터 공동 의사결정까지 다양한 수준으로 구현된다.

\begin{table}[htbp]
\centering
\caption{공공 AI 시민 참여 수준별 메커니즘}
\label{tab:18_citizen}
\begin{tabularx}{\textwidth}{p{2cm}p{3cm}Xp{2.5cm}}
\toprule
\textbf{참여 수준} & \textbf{메커니즘} & \textbf{내용} & \textbf{적용 단계} \\
\midrule
정보 공개 & 알고리즘 레지스터 & AI 시스템 목록 및 상세 정보 공개 & 운영 전 과정 \\
의견 수렴 & 공청회/의견 조회 & AI 도입 계획에 대한 시민 의견 수렴 & 기획/설계 \\
협의 & 시민 자문위원회 & AI 전문가 및 시민 대표 자문 & 설계/운영 \\
감시 & 시민 감사단 & AI 운영 결과 모니터링 및 평가 & 운영/평가 \\
공동 결정 & 참여 설계 워크숍 & AI 시스템 요구사항 공동 도출 & 기획 \\
\bottomrule
\end{tabularx}
\end{table}

\paragraph{공청회 운영 가이드}

공공 AI 도입 시 공청회는 시민의 우려를 사전에 파악하고 신뢰를 구축하는 핵심 절차이다. 공청회는 다음 요소를 포함해야 한다: (1) AI 시스템의 목적과 작동 방식을 비전문가 언어로 설명, (2) 예상되는 영향과 리스크 요소 공개, (3) 시민 질의에 대한 충분한 응답 시간, (4) 서면 의견 제출 기회, (5) 공청회 결과의 공식 반영 및 공개.

\paragraph{시민 감사단(Citizen Audit Panel)}

시민 감사단은 AI 시스템의 운영 결과를 시민의 관점에서 모니터링하는 조직이다. 감사단은 다양한 배경(연령, 지역, 직업, 장애 여부 등)의 시민으로 구성되며, 다음 역할을 수행한다:

\begin{itemize}
    \item AI 의사결정 결과의 형평성 모니터링(취약계층 불이익 여부)
    \item 시민 민원/불만 분석 및 패턴 보고
    \item AI 시스템 설명의 이해 가능성(Comprehensibility) 평가
    \item 개선 권고사항 도출 및 운영 기관에 전달
    \item 연례 시민 감사 보고서 작성 및 공개
\end{itemize}

\paragraph{알고리즘 자문위원회(Algorithm Advisory Board)}

알고리즘 자문위원회는 AI 전문가, 법률 전문가, 윤리학자, 시민 대표 등으로 구성되는 독립 자문 기구이다. 주요 역할은 다음과 같다:

\begin{table}[htbp]
\centering
\caption{알고리즘 자문위원회 구성 및 역할}
\label{tab:18_advisory}
\begin{tabularx}{\textwidth}{p{2.5cm}p{2cm}X}
\toprule
\textbf{구성원 유형} & \textbf{인원(예시)} & \textbf{역할} \\
\midrule
AI 기술 전문가 & 3명 & 알고리즘 적절성, 기술적 리스크 평가 \\
법률 전문가 & 2명 & 법적 준수 여부, 시민 권리 보호 검토 \\
윤리학자 & 1명 & 윤리적 쟁점 식별, 가치 충돌 조정 \\
시민 대표 & 3명 & 시민 관점 반영, 이해 가능성 평가 \\
운영 기관 담당자 & 2명 & 현장 운영 정보 제공, 개선 실행 \\
\bottomrule
\end{tabularx}
\end{table}

%----------------------------------------------------------
\section{공공 AI 조달과 검증}
\label{sec:18.3}

\subsection{AI 시스템 조달 가이드라인}
\label{sec:18.3.1}

공공 AI 조달은 단순 기능 구현이 아닌 성능 목표(정확도 등)와 윤리적 요건(공정성 등)을 중심으로 평가되어야 한다(표 \ref{tab:18_4}).

\begin{table}[htbp]
\centering
\caption{공공 AI 조달의 특수 고려사항}
\label{tab:18_4}
\begin{tabularx}{\textwidth}{p{3cm}Xp{5cm}}
\toprule
\textbf{고려사항} & \textbf{전통 IT 조달} & \textbf{AI 조달 특수성} \\
\midrule
요구사항 정의 & 기능 명세 가능 & 성능 목표(정확도, 공정성 등) 정의 \\
테스트 데이터 & 업체 제공 데이터 활용 & 발주기관 실데이터 기반 검증 필수 \\
유지보수 & 버그 수정 위주 & 모델 재학습, 드리프트 모니터링 포함 \\
윤리적 요건 & 명시적 규정 드묾 & 설명가능성, 편향 배제 요건 필수 \\
\bottomrule
\end{tabularx}
\end{table}

\paragraph{AI 조달 RFP 작성 가이드}

공공 AI 조달의 제안요청서(RFP, Request for Proposal) 작성 시에는 전통적 IT 조달과는 다른 특수 항목을 반드시 포함해야 한다.

\begin{table}[htbp]
\centering
\caption{공공 AI RFP 필수 포함 항목}
\label{tab:18_rfp}
\begin{tabularx}{\textwidth}{p{3cm}Xp{3.5cm}}
\toprule
\textbf{RFP 섹션} & \textbf{포함 내용} & \textbf{평가 기준} \\
\midrule
성능 요구사항 & 정확도, F1-Score, 지연시간 등 수치 목표 & 목표치 대비 달성률 \\
공정성 요구사항 & 보호 속성 명시, 4/5 Rule 준수 요건 & 편향 지표 충족 여부 \\
설명가능성 요구사항 & XAI 적용 의무, 시민 설명 템플릿 제공 & 설명 방식의 이해 가능성 \\
데이터 거버넌스 & 학습 데이터 출처/품질 관리 방안 & 데이터 관리 체계 적정성 \\
모니터링 요구사항 & 드리프트 탐지, 성능 모니터링 체계 & 모니터링 방안 구체성 \\
인간 개입 체계 & 자동화 수준, 에스컬레이션 절차 & 인간 제어권 보장 수준 \\
소스 코드 납품 & 소스 코드 및 모델 가중치 납품 범위 & 종속성 탈피(Vendor Lock-in 방지) \\
유지보수/재학습 & 모델 재학습 주기, 비용, 데이터 관리 & 운영 지속가능성 \\
보안 요구사항 & 적대적 공격 방어, 데이터 보호 & 보안 테스트 결과 \\
시민 참여 연계 & 시민 피드백 반영 메커니즘 & 시민 참여 체계 구비 \\
\bottomrule
\end{tabularx}
\end{table}

\paragraph{AI 조달 평가 기준}

공공 AI 조달 시 기술 평가 비중과 가격 평가 비중을 적절히 배분해야 한다. AI 시스템의 특성상 단순 가격 경쟁은 품질 저하로 이어질 수 있으므로, 기술 평가 비중을 70\% 이상으로 설정하는 것이 권장된다.

\begin{table}[htbp]
\centering
\caption{공공 AI 조달 평가 기준 배점표 (예시, 100점 만점)}
\label{tab:18_eval_criteria}
\begin{tabularx}{\textwidth}{p{3.5cm}p{1.5cm}X}
\toprule
\textbf{평가 항목} & \textbf{배점} & \textbf{평가 내용} \\
\midrule
기술 역량 & 15점 & 유사 프로젝트 수행 경험, AI 전문 인력 보유 \\
모델 성능 & 20점 & POC 결과, 정확도/F1/AUC 등 목표치 달성 \\
공정성/윤리 & 15점 & 편향 방지 전략, 공정성 테스트 계획 \\
설명가능성 & 10점 & XAI 구현 방안, 시민 설명 체계 \\
데이터 관리 & 10점 & 데이터 품질 관리, 가명처리 방안 \\
운영/모니터링 & 10점 & 드리프트 탐지, 성능 모니터링 계획 \\
보안 & 5점 & 적대적 공격 방어, 개인정보 보호 \\
가격 적정성 & 15점 & 총비용(TCO) 대비 가치 \\
\bottomrule
\end{tabularx}
\end{table}

\paragraph{AI 조달 계약 조건}

공공 AI 조달 계약에는 다음과 같은 특수 조건이 포함되어야 한다:

\begin{itemize}
    \item \textbf{성능 보증 조항}: 납품 후 6개월간 목표 성능(정확도, 공정성 지표) 유지 보증
    \item \textbf{모델 재학습 의무}: 성능 저하(PSI > 0.25 등) 발생 시 재학습 의무
    \item \textbf{소스 코드/모델 납품}: 벤더 종속(Vendor Lock-in) 방지를 위한 완전 납품
    \item \textbf{편향 시정 의무}: 운영 중 편향 발견 시 30일 이내 시정 의무
    \item \textbf{설명 체계 유지}: 시민 설명 요청 시 7일 이내 설명 제공 체계 유지
    \item \textbf{손해배상}: AI 오류로 인한 시민 피해 발생 시 업체 책임 범위 명시
\end{itemize}

\subsection{품질 인증 제도와 시민 참여}
\label{sec:18.3.2}

한국정보통신기술협회(TTA\cite{tta2024trust})의 GS 인증과 한국지능정보사회진흥원(NIA)의 AI 품질 인증은 공공 도입 시 핵심 지표가 된다 \cite{nia2024quality, tta2024trust}. 또한 민주적 정당성 확보를 위해 기획 단계의 공청회부터 운영 단계의 시민 모니터링단 운영까지 다각도의 \textbf{시민 참여(Citizen Participation)}가 요구된다.

\paragraph{NIA AI 품질 인증 상세}

한국지능정보사회진흥원(NIA)의 AI 품질 인증은 공공 AI 시스템의 품질을 객관적으로 검증하는 제도이다. 인증 체계는 기능 품질, 데이터 품질, 윤리 품질의 세 축으로 구성된다.

\begin{table}[htbp]
\centering
\caption{NIA AI 품질 인증 평가 체계}
\label{tab:18_nia}
\begin{tabularx}{\textwidth}{p{2.5cm}p{3cm}Xp{2cm}}
\toprule
\textbf{평가 영역} & \textbf{평가 항목} & \textbf{세부 기준} & \textbf{배점} \\
\midrule
기능 품질 & 정확성 & 목표 성능 지표 달성 여부 & 30점 \\
기능 품질 & 안정성 & 부하 테스트, 장애 복구 & 10점 \\
기능 품질 & 응답 시간 & 사용자 체감 응답 시간 & 10점 \\
데이터 품질 & 데이터 완전성 & 결측률, 이상치 관리 & 15점 \\
데이터 품질 & 데이터 최신성 & 학습 데이터 갱신 주기 & 5점 \\
윤리 품질 & 공정성 & 편향 지표 충족 여부 & 15점 \\
윤리 품질 & 설명가능성 & XAI 구현 수준 & 10점 \\
윤리 품질 & 개인정보 보호 & 가명처리/암호화 적정성 & 5점 \\
\bottomrule
\end{tabularx}
\end{table}

\paragraph{TTA GS 인증과 AI 신뢰성 인증}

TTA의 GS(Good Software) 인증은 전통적 소프트웨어 품질 인증이나, AI 시스템에 대해서는 별도의 \textbf{AI 신뢰성 인증}을 운영하고 있다. 이 인증은 ISO/IEC 25010(소프트웨어 품질 모델)과 ISO/IEC 42001(AI 경영시스템)을 기반으로 하며, 공정성, 견고성, 설명가능성 등 AI 고유 품질 속성을 평가한다.

%----------------------------------------------------------
\section{공공 AI 영향평가 자동화 도구}
\label{sec:18.3.3}

공공 AI 시스템의 도입과 운영 시 체계적인 영향평가를 자동화하는 것은 거버넌스 효율성과 일관성을 높이는 핵심 도구이다.

\begin{lstlisting}[language=Python, caption={Public AI Impact Assessment Automation Tool}, label={lst:public_impact}]
import json
import numpy as np
import pandas as pd
from dataclasses import dataclass, field
from typing import Dict, List, Optional, Tuple
from datetime import datetime
from enum import Enum

class RiskLevel(Enum):
    LOW = "low"
    MEDIUM = "medium"
    HIGH = "high"
    VERY_HIGH = "very_high"

class ImpactCategory(Enum):
    FUNDAMENTAL_RIGHTS = "fundamental_rights"
    EQUITY = "equity"
    PRIVACY = "privacy"
    TRANSPARENCY = "transparency"
    SAFETY = "safety"
    ACCOUNTABILITY = "accountability"

@dataclass
class ImpactAssessmentItem:
    category: ImpactCategory
    question: str
    score: int  # 1-5
    weight: float
    evidence: str
    mitigation: str

@dataclass
class PublicAIImpactReport:
    system_name: str
    agency: str
    assessment_date: datetime
    items: List[ImpactAssessmentItem]
    overall_risk: RiskLevel
    recommendations: List[str]
    approval_status: str

class PublicAIImpactAssessor:
    """Automated impact assessment tool for public AI systems"""

    ASSESSMENT_TEMPLATE = {
        ImpactCategory.FUNDAMENTAL_RIGHTS: [
            "AI system may restrict citizen's right to due process",
            "AI system may affect freedom of expression or assembly",
            "AI decisions may be difficult for citizens to challenge",
            "Vulnerable groups may be disproportionately affected",
        ],
        ImpactCategory.EQUITY: [
            "AI system produces different outcomes across regions",
            "Elderly or disabled citizens have reduced access",
            "Low-income groups face higher error rates",
            "Digital literacy gaps affect service quality",
        ],
        ImpactCategory.PRIVACY: [
            "Personal data is collected beyond minimum necessary",
            "Data retention exceeds required period",
            "Cross-agency data sharing lacks proper controls",
            "Re-identification risk from combined datasets",
        ],
        ImpactCategory.TRANSPARENCY: [
            "AI decision logic is not explainable to citizens",
            "System is not registered in algorithm register",
            "Performance metrics are not publicly available",
            "Citizens cannot request explanation of decisions",
        ],
        ImpactCategory.SAFETY: [
            "System failure could cause significant citizen harm",
            "No fallback procedure exists for system outage",
            "Adversarial attacks could manipulate outcomes",
            "No regular security assessment is performed",
        ],
        ImpactCategory.ACCOUNTABILITY: [
            "Responsibility chain for AI errors is unclear",
            "No human oversight mechanism exists",
            "Appeal/grievance process is not established",
            "No regular audit cycle is defined",
        ],
    }

    CATEGORY_WEIGHTS = {
        ImpactCategory.FUNDAMENTAL_RIGHTS: 0.25,
        ImpactCategory.EQUITY: 0.20,
        ImpactCategory.PRIVACY: 0.15,
        ImpactCategory.TRANSPARENCY: 0.15,
        ImpactCategory.SAFETY: 0.15,
        ImpactCategory.ACCOUNTABILITY: 0.10,
    }

    def __init__(self):
        self.assessments: List[PublicAIImpactReport] = []

    def conduct_assessment(self, system_name: str, agency: str,
                           responses: Dict[str, List[Dict]]) -> PublicAIImpactReport:
        """Conduct automated impact assessment"""
        items = []

        for category, questions in self.ASSESSMENT_TEMPLATE.items():
            cat_responses = responses.get(category.value, [])
            weight = self.CATEGORY_WEIGHTS[category]

            for i, question in enumerate(questions):
                resp = cat_responses[i] if i < len(cat_responses) else {}
                items.append(ImpactAssessmentItem(
                    category=category,
                    question=question,
                    score=resp.get("score", 3),
                    weight=weight / len(questions),
                    evidence=resp.get("evidence", ""),
                    mitigation=resp.get("mitigation", "")
                ))

        # Calculate overall risk
        weighted_score = sum(item.score * item.weight for item in items)
        if weighted_score >= 4.0:
            overall_risk = RiskLevel.VERY_HIGH
        elif weighted_score >= 3.0:
            overall_risk = RiskLevel.HIGH
        elif weighted_score >= 2.0:
            overall_risk = RiskLevel.MEDIUM
        else:
            overall_risk = RiskLevel.LOW

        # Generate recommendations
        recommendations = self._generate_recommendations(items, overall_risk)

        # Determine approval status
        if overall_risk == RiskLevel.VERY_HIGH:
            approval = "REJECTED - Major redesign required"
        elif overall_risk == RiskLevel.HIGH:
            approval = "CONDITIONAL - Mitigation measures required before deployment"
        elif overall_risk == RiskLevel.MEDIUM:
            approval = "APPROVED WITH CONDITIONS - Enhanced monitoring required"
        else:
            approval = "APPROVED - Standard monitoring"

        report = PublicAIImpactReport(
            system_name=system_name,
            agency=agency,
            assessment_date=datetime.now(),
            items=items,
            overall_risk=overall_risk,
            recommendations=recommendations,
            approval_status=approval
        )

        self.assessments.append(report)
        return report

    def _generate_recommendations(self, items: List[ImpactAssessmentItem],
                                   risk: RiskLevel) -> List[str]:
        """Generate context-specific recommendations"""
        recs = []
        high_risk_items = [i for i in items if i.score >= 4]

        for item in high_risk_items:
            if item.category == ImpactCategory.FUNDAMENTAL_RIGHTS:
                recs.append("Implement mandatory human review for high-impact decisions")
            elif item.category == ImpactCategory.EQUITY:
                recs.append("Conduct quarterly equity audit with disaggregated data")
            elif item.category == ImpactCategory.PRIVACY:
                recs.append("Strengthen data minimization and retention policies")
            elif item.category == ImpactCategory.TRANSPARENCY:
                recs.append("Register system in algorithm transparency register")
            elif item.category == ImpactCategory.SAFETY:
                recs.append("Develop and test fallback procedures")
            elif item.category == ImpactCategory.ACCOUNTABILITY:
                recs.append("Establish clear accountability chain and audit schedule")

        if risk in [RiskLevel.HIGH, RiskLevel.VERY_HIGH]:
            recs.append("Conduct public consultation before deployment")
            recs.append("Establish citizen audit panel for ongoing monitoring")

        return recs

    def export_report(self, report: PublicAIImpactReport) -> Dict:
        """Export assessment report in structured format"""
        return {
            "system_name": report.system_name,
            "agency": report.agency,
            "assessment_date": report.assessment_date.isoformat(),
            "overall_risk": report.overall_risk.value,
            "approval_status": report.approval_status,
            "category_scores": {
                cat.value: np.mean([i.score for i in report.items if i.category == cat])
                for cat in ImpactCategory
            },
            "recommendations": report.recommendations,
            "total_items_assessed": len(report.items),
            "high_risk_items": sum(1 for i in report.items if i.score >= 4)
        }
\end{lstlisting}

%----------------------------------------------------------
\section{구현 사례 연구}
\label{sec:18.4}

\subsection{행정 자동화: 국민연금공단 AI 기반 자격심사}

\begin{practicaltool}{국민연금공단 AI 기반 자격심사 시스템}
연간 약 50만 건의 노령연금 심사를 자동화하며 구축한 거버넌스 모델이다.\\
\textbf{자동화 범위}: 명확한 수급 요건(기속 행위) 중심 자동화, 특수 경력 등 재량 영역은 인간 검토.\\
\textbf{공정성}: 분기별 지역/연령별 승인율 편차 분석, 10\% 이상 격차 시 모델 전수 조사.\\
\textbf{투명성}: 신청자에게 결과와 함께 결정 근거 3가지를 제시하고 불복 절차 상시 안내.
\end{practicaltool}

\subsection{교육 AI 거버넌스 사례}

교육 분야의 공공 AI는 학생의 학습 성과, 진로 추천, 학교 배정 등에 활용되며, 특히 미성년자에 대한 AI 의사결정이라는 점에서 더욱 엄격한 거버넌스가 요구된다.

\begin{practicaltool}{교육부 AI 디지털교과서 거버넌스 사례}
2025년부터 도입되는 AI 디지털교과서의 학습 분석(Learning Analytics) 기능에 대한 거버넌스 체계이다.\\
\textbf{데이터 보호}: 학생 학습 데이터의 수집 범위를 학습 활동에 엄격히 한정하며, 행동 패턴 수집을 금지한다. 데이터 보존 기간은 해당 학년도 종료 후 1년으로 제한한다.\\
\textbf{공정성}: AI 추천 학습 콘텐츠가 특정 성별, 지역, 소득 수준의 학생에게 불리하지 않도록 분기별 형평성 평가를 수행한다.\\
\textbf{투명성}: 학생과 학부모에게 AI가 학습 수준을 어떻게 판단했는지, 왜 특정 콘텐츠를 추천했는지를 이해 가능한 언어로 설명한다.\\
\textbf{인간 개입}: AI의 학습 수준 진단은 참고 자료로만 활용하며, 최종 평가와 진로 지도는 반드시 교사가 수행한다.\\
\textbf{보호자 동의}: 18세 미만 학생의 AI 학습 데이터 활용 시 보호자 동의를 필수로 받으며, 옵트아웃(Opt-out) 시 비AI 대안을 제공한다.
\end{practicaltool}

\begin{table}[htbp]
\centering
\caption{교육 AI 거버넌스 핵심 원칙}
\label{tab:18_edu_ai}
\begin{tabularx}{\textwidth}{p{2.5cm}Xp{4cm}}
\toprule
\textbf{원칙} & \textbf{내용} & \textbf{구현 방안} \\
\midrule
아동 최선의 이익 & AI가 학생의 성장과 발달에 기여 & 학습 격차 해소 목표 우선 설정 \\
데이터 최소 수집 & 학습 목적 외 데이터 수집 금지 & 수집 항목 사전 심사, 자동 삭제 \\
비분류 원칙 & 학생을 고정 등급으로 분류 금지 & 성장 가능성 기반 동적 평가 \\
교사 주도성 & AI는 도구, 최종 판단은 교사 & 교사 대시보드, AI 제안 수용/거부권 \\
보호자 참여 & 보호자의 알 권리와 통제권 & 정기 보고서, 옵트아웃 보장 \\
\bottomrule
\end{tabularx}
\end{table}

\subsection{민원상담 챗봇 거버넌스 상세}

민원상담 챗봇은 공공기관에서 가장 빠르게 확산되는 AI 응용 영역이다. 그러나 챗봇의 답변 오류가 시민의 권리 행사에 직접 영향을 미칠 수 있어, 체계적인 거버넌스가 필수적이다.

\begin{practicaltool}{정부24 민원상담 챗봇 거버넌스}
정부24의 민원상담 챗봇은 연간 약 200만 건의 상담을 처리하며 다음과 같은 거버넌스 체계를 구축하였다.\\
\textbf{답변 범위 제한}: 챗봇은 사실 정보 안내(접수처, 구비 서류, 절차 안내)에 한정하며, 법적 해석이나 판단이 필요한 질문은 즉시 담당 공무원에게 에스컬레이션한다.\\
\textbf{답변 정확성}: 챗봇의 답변은 법령 데이터베이스와 실시간 연동되며, 법령 변경 시 24시간 이내에 답변 콘텐츠가 자동 업데이트된다.\\
\textbf{오답 관리}: 시민이 챗봇 답변에 대해 '부정확' 피드백을 제공할 수 있으며, 동일 질문에 대한 부정확 피드백이 3건 이상 누적되면 해당 답변을 자동 비활성화하고 인간 검토를 요청한다.\\
\textbf{접근성}: 음성 인식 기반 상담(시각장애인), 쉬운 한국어 모드(외국인, 노인), 수어 안내 연계(청각장애인)를 제공한다.\\
\textbf{개인정보 보호}: 상담 내용은 서비스 개선 목적으로만 6개월간 보관하며, 개인 식별 정보는 상담 종료 즉시 삭제한다.
\end{practicaltool}

\begin{lstlisting}[language=Python, caption={Public Chatbot Governance: Answer Quality and Escalation Management}, label={lst:chatbot_gov}]
from dataclasses import dataclass, field
from typing import Dict, List, Optional
from datetime import datetime, timedelta
from enum import Enum
from collections import defaultdict

class QueryCategory(Enum):
    FACTUAL_INFO = "factual_info"        # (automatable)
    PROCEDURE_GUIDE = "procedure_guide"  # (automatable)
    LEGAL_INTERPRETATION = "legal"       # (escalation required)
    COMPLAINT = "complaint"              # (escalation required)
    COMPLEX_CASE = "complex"             # (escalation required)

class AnswerStatus(Enum):
    ACTIVE = "active"
    UNDER_REVIEW = "under_review"
    DEACTIVATED = "deactivated"
    UPDATED = "updated"

@dataclass
class ChatbotAnswer:
    answer_id: str
    question_pattern: str
    answer_text: str
    category: QueryCategory
    status: AnswerStatus
    legal_reference: str
    last_verified: datetime
    negative_feedback_count: int = 0
    total_served: int = 0

class PublicChatbotGovernance:
    """Governance system for public service chatbot"""

    ESCALATION_CATEGORIES = [
        QueryCategory.LEGAL_INTERPRETATION,
        QueryCategory.COMPLAINT,
        QueryCategory.COMPLEX_CASE
    ]

    def __init__(self, negative_feedback_threshold: int = 3,
                 answer_refresh_days: int = 30):
        self.answers: Dict[str, ChatbotAnswer] = {}
        self.neg_threshold = negative_feedback_threshold
        self.refresh_days = answer_refresh_days
        self.escalation_log: List[Dict] = []

    def process_query(self, query_text: str, category: QueryCategory,
                      matched_answer_id: Optional[str]) -> Dict:
        """Process citizen query with governance checks"""

        # Check if escalation is required
        if category in self.ESCALATION_CATEGORIES:
            self._log_escalation(query_text, category, "Category requires human review")
            return {
                "action": "escalate",
                "message": "This question requires expert review. "
                          "Transferring to a staff member.",
                "estimated_response": "Within 2 business hours"
            }

        if not matched_answer_id or matched_answer_id not in self.answers:
            self._log_escalation(query_text, category, "No matching answer found")
            return {
                "action": "escalate",
                "message": "Unable to find a matching answer. "
                          "Transferring to a staff member."
            }

        answer = self.answers[matched_answer_id]

        # Check answer status
        if answer.status != AnswerStatus.ACTIVE:
            self._log_escalation(query_text, category,
                               f"Answer status: {answer.status.value}")
            return {"action": "escalate", "message": "Answer is under review."}

        # Check answer freshness
        if (datetime.now() - answer.last_verified).days > self.refresh_days:
            answer.status = AnswerStatus.UNDER_REVIEW
            return {"action": "escalate",
                    "message": "Answer requires verification update."}

        answer.total_served += 1
        return {
            "action": "respond",
            "answer": answer.answer_text,
            "legal_reference": answer.legal_reference,
            "feedback_prompt": "Was this answer helpful? [Yes/No]"
        }

    def record_feedback(self, answer_id: str, is_positive: bool):
        """Record citizen feedback on answer quality"""
        if answer_id not in self.answers:
            return

        answer = self.answers[answer_id]
        if not is_positive:
            answer.negative_feedback_count += 1

            if answer.negative_feedback_count >= self.neg_threshold:
                answer.status = AnswerStatus.DEACTIVATED
                self._log_escalation(
                    f"Auto-deactivated: {answer.question_pattern}",
                    QueryCategory.FACTUAL_INFO,
                    f"Negative feedback count: {answer.negative_feedback_count}"
                )

    def _log_escalation(self, query: str, category: QueryCategory, reason: str):
        """Log escalation for audit trail"""
        self.escalation_log.append({
            "timestamp": datetime.now().isoformat(),
            "query": query[:200],
            "category": category.value,
            "reason": reason
        })

    def generate_quality_report(self) -> Dict:
        """Generate monthly quality report"""
        total_answers = len(self.answers)
        active = sum(1 for a in self.answers.values()
                     if a.status == AnswerStatus.ACTIVE)
        deactivated = sum(1 for a in self.answers.values()
                         if a.status == AnswerStatus.DEACTIVATED)
        total_escalations = len(self.escalation_log)

        return {
            "report_period": datetime.now().strftime("%Y-%m"),
            "total_answers_in_db": total_answers,
            "active_answers": active,
            "deactivated_answers": deactivated,
            "total_escalations": total_escalations,
            "answer_health_rate": active / total_answers if total_answers > 0 else 0
        }
\end{lstlisting}

\subsection{한국도로공사 지능형 교통 시스템}

\begin{practicaltool}{고속도로 AI 교통관리 시스템 거버넌스}
전국 고속도로망의 실시간 돌발상황을 탐지하고 관제하는 초연결 거버넌스 사례이다.\\
\textbf{Human-in-the-Loop}: AI의 모든 제어 권고(속도 제한, 차로 제어)는 관제사 최종 승인 후 실행.\\
\textbf{안전 무결성}: 사고, 역주행 등 고위험 상황은 AI 자동 경보와 동시에 관제사 즉시 개입 유도.\\
\textbf{프라이버시}: CCTV 영상의 차량 번호/얼굴 실시간 비식별화 및 보존 기간 엄격 제한.\\
\textbf{기술적 고도화}: 2026년 기준 엣지 AI를 활용하여 현장에서 사고 30초 내 탐지 및 자동 알림 체계 구축.
\end{practicaltool}

\begin{lstlisting}[language=Python, caption={ITS AI Surveillance and Safety Governance Management System}, label={lst:its_gov}]
class ITSGovernanceManager:
 """Proper Role Allocation between ITS AI Recommendations and Human Control"""
 def __init__(self):
 self.auto_execute_types = [IncidentType.CONGESTION] # self.critical_types = [IncidentType.ACCIDENT, IncidentType.WRONG_WAY]

 def process_recommendation(self, rec):
 # (Safety-First)
 if rec.incident_type in self.critical_types:
 return {"action": "human_review_required", "urgency": "EMERGENCY"}

 # 95% days (Efficiency)
 if rec.incident_type in self.auto_execute_types and rec.confidence > 0.95:
 return {"action": "auto_execute", "logging": True}

 return {"action": "human_review_required", "urgency": "NORMAL"}
\end{lstlisting}

%----------------------------------------------------------
\section{공공 AI 조달 RFP 템플릿}
\label{sec:18.5}

\begin{practicaltool}{공공 AI 조달 RFP 템플릿}
\textbf{1. 사업 개요}
\begin{itemize}
    \item 사업명: [OO기관 AI 기반 OO 시스템 구축]
    \item 사업 기간: [YYYY.MM -- YYYY.MM]
    \item 사업 예산: [OO백만원]
    \item AI 시스템 목적: [구체적 행정 업무 및 자동화 범위 기술]
    \item 대상 시민: [서비스 수혜 대상 및 예상 이용 규모]
\end{itemize}

\textbf{2. 성능 요구사항}
\begin{itemize}
    \item 정확도: [목표 지표(예: F1-Score > 0.90)]
    \item 응답 시간: [실시간 요건(예: 95\%ile < 500ms)]
    \item 가용성: [SLA(예: 99.9\% 가용률)]
    \item 처리 용량: [동시 처리 건수(예: 1,000 TPS)]
\end{itemize}

\textbf{3. 윤리/공정성 요구사항}
\begin{itemize}
    \item 보호 속성: [성별, 연령, 지역, 장애 여부 등 명시]
    \item 공정성 지표: [4/5 Rule(Disparate Impact > 0.8) 필수 충족]
    \item 설명가능성: [XAI 도구 적용 필수, 시민 설명 템플릿 제공]
    \item 영향평가: [도입 전 영향평가 수행 및 결과 제출]
\end{itemize}

\textbf{4. 데이터 거버넌스 요구사항}
\begin{itemize}
    \item 학습 데이터: [발주기관 제공 데이터 기반 학습 필수]
    \item 가명처리: [개인정보보호법 준수 가명처리 방안 제시]
    \item 데이터 보존: [학습 완료 후 원본 데이터 파기 증명]
    \item 데이터 품질: [결측률 < 5\%, 이상치 처리 기준 명시]
\end{itemize}

\textbf{5. 운영/유지보수 요구사항}
\begin{itemize}
    \item 모니터링: [성능 모니터링 대시보드 구축, PSI 기반 드리프트 탐지]
    \item 재학습: [분기별 재학습 계획 및 비용 포함]
    \item 비상 절차: [AI 장애 시 수동 전환 절차 수립]
    \item 인간 개입: [에스컬레이션 기준 및 절차 명시]
\end{itemize}

\textbf{6. 납품/인수 요구사항}
\begin{itemize}
    \item 소스 코드: [전체 소스 코드 및 모델 가중치 납품]
    \item 문서: [시스템 설계서, 모델 개발 문서, 운영 매뉴얼]
    \item 교육: [운영 담당자 교육 (최소 40시간)]
    \item 인증: [NIA AI 품질 인증 또는 TTA AI 신뢰성 인증 취득]
\end{itemize}

\textbf{7. 시민 참여 연계}
\begin{itemize}
    \item 알고리즘 레지스터: [투명성 레지스터 등록 정보 제공]
    \item 시민 피드백: [시민 피드백 수집 및 반영 메커니즘 구축]
    \item 이의제기: [시민 이의제기 접수/처리 체계 구축]
\end{itemize}
\end{practicaltool}

%----------------------------------------------------------
\section{시민 참여 운영 가이드}
\label{sec:18.6}

\begin{practicaltool}{공공 AI 시민 참여 운영 가이드}
\textbf{1. 기획 단계 시민 참여}
\begin{itemize}
    \item AI 도입 목적과 범위를 비전문가 언어로 설명하는 공청회 개최
    \item 공청회 개최 최소 30일 전 사전 안내 및 자료 공개
    \item 서면 의견 제출 기간 최소 14일 보장
    \item 수렴된 의견의 반영/미반영 사유를 공식 문서로 회신
\end{itemize}

\textbf{2. 설계 단계 시민 참여}
\begin{itemize}
    \item 알고리즘 자문위원회 구성 (AI 전문가 + 법률 전문가 + 시민 대표)
    \item 자문위원회 분기별 회의, 회의록 공개
    \item 취약계층 대표자를 반드시 포함 (장애인, 노인, 외국인 등)
    \item 설계 대안에 대한 시민 선호도 조사 실시
\end{itemize}

\textbf{3. 운영 단계 시민 참여}
\begin{itemize}
    \item 시민 감사단 구성 (10-20명, 다양한 배경)
    \item 감사단 월 1회 운영 결과 리뷰 회의
    \item 시민 만족도 조사 분기별 실시, 결과 공개
    \item 민원/불만 분석 리포트 월간 생성 및 개선 반영
\end{itemize}

\textbf{4. 평가 단계 시민 참여}
\begin{itemize}
    \item 연례 시민 감사 보고서 발간 및 공개
    \item AI 시스템 지속/변경/폐기 결정에 시민 의견 반영
    \item 형평성 평가 결과를 시민이 이해할 수 있는 형태로 공개
    \item 차기 연도 개선 계획에 시민 권고사항 반영
\end{itemize}
\end{practicaltool}

%----------------------------------------------------------
\section{공공 AI 거버넌스 성숙도 모델}
\label{sec:18.7}

공공기관이 AI 거버넌스를 체계적으로 발전시키기 위해서는 현재 수준을 진단하고 목표를 설정할 수 있는 성숙도 모델이 필요하다.

\begin{table}[htbp]
\centering
\caption{공공 AI 거버넌스 성숙도 모델}
\label{tab:18_maturity}
\begin{tabularx}{\textwidth}{p{1.5cm}p{2cm}Xp{4cm}}
\toprule
\textbf{수준} & \textbf{단계} & \textbf{특성} & \textbf{주요 활동} \\
\midrule
1 & 초기 & AI 도입이 산발적, 거버넌스 부재 & 개별 부서별 AI 시범 사업 \\
2 & 인식 & 거버넌스 필요성 인식, 기본 정책 수립 & AI 윤리 원칙 선언, 담당 조직 지정 \\
3 & 체계화 & 표준 프로세스 수립, 조직 구축 & 영향평가 절차, 검증 체계, 시민 참여 \\
4 & 관리 & 정량적 관리, 지속적 모니터링 & KPI 관리, 대시보드, 정기 감사 \\
5 & 최적화 & 지속적 개선, 선도적 거버넌스 & 혁신적 참여 모델, 글로벌 표준 주도 \\
\bottomrule
\end{tabularx}
\end{table}

\subsection{성숙도 수준별 핵심 역량}

각 성숙도 수준에서 갖추어야 할 핵심 역량은 다음과 같다.

\begin{table}[htbp]
\centering
\caption{공공 AI 거버넌스 성숙도 수준별 역량 매트릭스}
\label{tab:18_maturity_detail}
\begin{tabularx}{\textwidth}{p{2cm}p{2.5cm}p{2.5cm}p{2.5cm}X}
\toprule
\textbf{영역} & \textbf{수준 2} & \textbf{수준 3} & \textbf{수준 4} & \textbf{수준 5} \\
\midrule
정책 & AI 윤리 원칙 & 거버넌스 규정 & 세부 가이드라인 & 자율 규제 체계 \\
조직 & 담당자 지정 & 전담 조직 & CDO/CAIO 설치 & 범정부 거버넌스 \\
프로세스 & 기본 검토 & 영향평가 수행 & 표준 운영 절차 & 자동화된 점검 \\
기술 & XAI 시범 & 모니터링 구축 & 실시간 대시보드 & AI 자가 진단 \\
시민 참여 & 정보 공개 & 의견 수렴 & 시민 감사단 & 공동 의사결정 \\
\bottomrule
\end{tabularx}
\end{table}

\subsection{성숙도 자가 진단 도구}

공공기관이 현재 AI 거버넌스 성숙도를 자가 진단하기 위한 체크리스트를 제공한다.

\begin{lstlisting}[language=Python, caption={Public AI Governance Maturity Self-Assessment Tool}, label={lst:maturity}]
from dataclasses import dataclass
from typing import Dict, List
from enum import Enum

class MaturityLevel(Enum):
    INITIAL = 1
    AWARE = 2
    SYSTEMATIC = 3
    MANAGED = 4
    OPTIMIZED = 5

class AssessmentDomain(Enum):
    POLICY = "policy"
    ORGANIZATION = "organization"
    PROCESS = "process"
    TECHNOLOGY = "technology"
    CITIZEN_PARTICIPATION = "citizen_participation"
    DATA_GOVERNANCE = "data_governance"
    RISK_MANAGEMENT = "risk_management"

@dataclass
class DomainAssessment:
    domain: AssessmentDomain
    score: int  # 1-5
    evidence: str
    improvement_actions: List[str]

class PublicAIMaturityAssessor:
    """Self-assessment tool for public AI governance maturity"""

    DOMAIN_CRITERIA = {
        AssessmentDomain.POLICY: {
            1: "No AI governance policy exists",
            2: "Basic AI ethics principles declared",
            3: "Formal governance regulations established",
            4: "Detailed guidelines per AI system type",
            5: "Self-regulating adaptive governance framework"
        },
        AssessmentDomain.ORGANIZATION: {
            1: "No designated AI governance personnel",
            2: "Part-time AI governance coordinator assigned",
            3: "Dedicated AI governance team established",
            4: "CDO/CAIO role with authority and budget",
            5: "Cross-agency governance coordination body"
        },
        AssessmentDomain.PROCESS: {
            1: "Ad-hoc AI deployment without review",
            2: "Basic review before AI deployment",
            3: "Standardized impact assessment process",
            4: "Automated compliance checking integrated",
            5: "Continuous improvement with feedback loops"
        },
        AssessmentDomain.TECHNOLOGY: {
            1: "No monitoring or XAI tools",
            2: "Basic performance logging",
            3: "XAI and bias monitoring deployed",
            4: "Real-time governance dashboard",
            5: "AI self-diagnosis and auto-remediation"
        },
        AssessmentDomain.CITIZEN_PARTICIPATION: {
            1: "No citizen engagement mechanism",
            2: "Information disclosure only",
            3: "Public hearings and comment periods",
            4: "Active citizen audit panels",
            5: "Co-design and participatory governance"
        },
        AssessmentDomain.DATA_GOVERNANCE: {
            1: "No data governance for AI",
            2: "Basic data quality checks",
            3: "Pseudonymization and access controls",
            4: "Data fabric with automated compliance",
            5: "Privacy-by-design with citizen data sovereignty"
        },
        AssessmentDomain.RISK_MANAGEMENT: {
            1: "No AI risk assessment",
            2: "Initial risk identification",
            3: "Structured risk assessment framework",
            4: "Quantitative risk monitoring with KPIs",
            5: "Predictive risk management with early warning"
        }
    }

    def __init__(self):
        self.assessments: List[DomainAssessment] = []

    def assess_domain(self, domain: AssessmentDomain,
                      score: int, evidence: str) -> DomainAssessment:
        """Assess a single governance domain"""
        improvements = []
        if score < 5:
            next_level = self.DOMAIN_CRITERIA[domain].get(score + 1, "")
            improvements.append(f"Target: {next_level}")

        assessment = DomainAssessment(
            domain=domain, score=score,
            evidence=evidence, improvement_actions=improvements
        )
        self.assessments.append(assessment)
        return assessment

    def calculate_overall_maturity(self) -> Dict:
        """Calculate overall maturity level"""
        if not self.assessments:
            return {"level": MaturityLevel.INITIAL, "score": 1.0}

        avg_score = sum(a.score for a in self.assessments) / len(self.assessments)

        # Overall level is constrained by the lowest domain
        min_score = min(a.score for a in self.assessments)
        effective_level = min(int(avg_score), min_score + 1)

        level = MaturityLevel(min(effective_level, 5))

        return {
            "level": level,
            "average_score": round(avg_score, 2),
            "min_domain_score": min_score,
            "weakest_domain": min(self.assessments, key=lambda a: a.score).domain.value,
            "domain_scores": {a.domain.value: a.score for a in self.assessments},
            "recommendation": f"Focus on improving '{min(self.assessments, key=lambda a: a.score).domain.value}' "
                            f"to advance overall maturity"
        }
\end{lstlisting}

%----------------------------------------------------------
\section{복지 AI 거버넌스 상세}
\label{sec:18.8}

복지 분야의 공공 AI는 수급자 선정, 급여 산정, 부정 수급 탐지 등에 활용되며, 시민의 생존권과 직결된다는 점에서 가장 높은 수준의 거버넌스가 요구된다.

\begin{table}[htbp]
\centering
\caption{복지 AI 유형별 거버넌스 요건}
\label{tab:18_welfare}
\begin{tabularx}{\textwidth}{p{2.5cm}p{2.5cm}Xp{3cm}}
\toprule
\textbf{AI 유형} & \textbf{리스크 수준} & \textbf{핵심 거버넌스 요건} & \textbf{자동화 허용 범위} \\
\midrule
수급 자격 판정 & 매우 높음 & 오판 시 생존권 침해, 100\% 설명 필수 & 기속 행위만 자동화 \\
급여액 산정 & 높음 & 산정 기준 투명 공개, 이의제기 보장 & 반자동화(AI 보조) \\
부정 수급 탐지 & 높음 & 차별적 감시 방지, 무죄 추정 & AI 제안, 인간 최종 결정 \\
맞춤형 서비스 추천 & 중간 & 정보 비대칭 방지, 옵트아웃 보장 & 자동화 가능 (정보 제공) \\
\bottomrule
\end{tabularx}
\end{table}

\paragraph{부정 수급 탐지 AI의 특수 거버넌스}

부정 수급 탐지 AI는 네덜란드 SyRI와 호주 Robodebt 사건에서 보았듯이 가장 높은 거버넌스 리스크를 내포한다. 핵심 원칙은 다음과 같다:

\begin{itemize}
    \item \textbf{무죄 추정 원칙}: AI의 이상 탐지 결과는 '조사 시작의 근거'일 뿐, 부정 수급의 증거가 아니다.
    \item \textbf{차별적 감시 금지}: 특정 지역, 소득 수준, 국적에 대한 과도한 감시를 방지하기 위해 탐지 대상의 인구통계학적 분포를 정기적으로 점검한다.
    \item \textbf{인간 최종 결정}: AI 탐지 결과를 바탕으로 한 조치(급여 중단, 환수 통보 등)는 반드시 담당 공무원이 사안을 개별 검토한 후 결정한다.
    \item \textbf{이의제기 보장}: 부정 수급으로 판단된 시민에게 AI 탐지 근거를 설명하고, 이의를 제기할 수 있는 실질적 절차를 제공한다.
    \item \textbf{정기 감사}: 부정 수급 탐지 AI의 정확도(정밀도, 재현율)와 형평성 지표를 분기별로 감사한다.
\end{itemize}

\paragraph{복지 사각지대 발굴 AI}

복지 사각지대 발굴 AI는 공공 AI의 긍정적 활용 사례로, 복지 혜택을 받을 자격이 있으나 신청하지 않은 시민을 사전에 발굴하여 안내하는 시스템이다. 이 AI는 다양한 행정 데이터(건강보험 자격 변동, 국세 소득 변동, 주민등록 변동 등)를 결합하여 위기 가구를 식별한다.

\begin{table}[htbp]
\centering
\caption{복지 사각지대 발굴 AI 거버넌스 체계}
\label{tab:18_welfare_gap}
\begin{tabularx}{\textwidth}{p{2.5cm}Xp{4cm}}
\toprule
\textbf{거버넌스 항목} & \textbf{내용} & \textbf{기준/절차} \\
\midrule
데이터 결합 & 8개 기관 33종 데이터 결합 & 결합 전문기관 경유, 가명처리 \\
모델 정확도 & 실제 위기 가구 발굴 비율 & 재현율 > 70\% 목표 \\
오탐 관리 & 비위기 가구를 위기로 판정 & 오탐 시 안내 문자만 발송, 불이익 없음 \\
프라이버시 & 시민 감시 우려 방지 & 발굴 결과 열람 권한 엄격 제한 \\
옵트아웃 & 발굴 대상 제외 요청 & 시민 요청 시 즉시 제외 \\
\bottomrule
\end{tabularx}
\end{table}

%----------------------------------------------------------
\section{공공 AI 거버넌스 인프라 구축}
\label{sec:18.9}

\subsection{공공 AI 거버넌스 플랫폼 아키텍처}

공공기관 전체의 AI 거버넌스를 통합 관리하기 위해서는 거버넌스 플랫폼이 필요하다. 이 플랫폼은 알고리즘 레지스터, 영향평가 시스템, 모니터링 대시보드, 시민 피드백 시스템을 통합한다.

\begin{table}[htbp]
\centering
\caption{공공 AI 거버넌스 플랫폼 구성 요소}
\label{tab:18_platform}
\begin{tabularx}{\textwidth}{p{3cm}Xp{3.5cm}}
\toprule
\textbf{구성 요소} & \textbf{기능} & \textbf{연동 대상} \\
\midrule
AI 시스템 레지스터 & 전체 AI 시스템 등록/관리 & 각 기관 AI 시스템 \\
영향평가 시스템 & 온라인 영향평가 수행/관리 & 알고리즘 자문위원회 \\
모니터링 대시보드 & 성능/공정성/안정성 실시간 감시 & AI 시스템 API \\
시민 포털 & 레지스터 열람, 피드백, 이의제기 & 정부24, 시민 인증 \\
감사 관리 & 감사 일정, 결과, 개선 추적 & 감사 조직 \\
보고서 생성기 & 국회/의회 보고, 공개 보고서 자동 생성 & 전체 구성 요소 \\
\bottomrule
\end{tabularx}
\end{table}

\subsection{공공 AI 인력 역량 강화}

공공 AI 거버넌스의 실효성은 궁극적으로 이를 운영하는 공무원의 역량에 달려 있다. AI에 대한 이해가 부족한 상태에서는 형식적인 거버넌스에 그칠 위험이 있다.

\begin{table}[htbp]
\centering
\caption{공공 AI 거버넌스 인력 역량 체계}
\label{tab:18_capacity}
\begin{tabularx}{\textwidth}{p{2.5cm}p{3cm}Xp{2.5cm}}
\toprule
\textbf{역할} & \textbf{필요 역량} & \textbf{교육 내용} & \textbf{교육 시간} \\
\midrule
정책 결정자 & AI 리터러시 & AI 기본 개념, 리스크/기회, 사례 & 16시간 \\
사업 관리자 & AI 프로젝트 관리 & AI 조달, 검증, 모니터링 방법론 & 40시간 \\
기술 담당자 & AI 기술 운영 & 모델 평가, 편향 테스트, XAI 활용 & 80시간 \\
시민 소통 담당 & AI 설명 역량 & 비전문가 언어 설명법, 민원 대응 & 24시간 \\
감사 담당자 & AI 감사 역량 & AI 감사 방법론, 증빙 검증 & 40시간 \\
\bottomrule
\end{tabularx}
\end{table}

\subsection{공공 AI 윤리 교육 체계}

모든 공공 AI 관련 종사자는 정기적인 윤리 교육을 이수해야 한다. 교육 체계는 기본 교육(AI 윤리 원칙, 공공성 가치), 심화 교육(편향 사례 분석, 영향평가 실습), 전문 교육(기술적 공정성 테스트, XAI 구현)의 3단계로 구성된다.

\begin{practicaltool}{공공 AI 거버넌스 종합 점검표}
\textbf{법적 준수}
\begin{itemize}
    \item 행정기본법 제20조 자동적 처분 요건 충족 여부 \checkmark / \texttimes
    \item 지능정보화 기본법 안전성 평가 수행 여부 \checkmark / \texttimes
    \item 개인정보보호법 가명처리/결합 절차 준수 \checkmark / \texttimes
    \item AI 기본법 고위험 AI 영향평가 수행 여부 \checkmark / \texttimes
\end{itemize}

\textbf{시민 권리 보장}
\begin{itemize}
    \item 알고리즘 투명성 레지스터 등록 여부 \checkmark / \texttimes
    \item 시민 설명 요청권 보장 체계 구축 \checkmark / \texttimes
    \item 이의제기/불복 절차 상시 안내 \checkmark / \texttimes
    \item 오프라인 대안 서비스 제공 \checkmark / \texttimes
\end{itemize}

\textbf{형평성 보장}
\begin{itemize}
    \item 취약계층 영향 평가 수행 여부 \checkmark / \texttimes
    \item 지역/연령/장애 등 집단별 편차 모니터링 \checkmark / \texttimes
    \item 디지털 접근성(WCAG 2.1 AA) 준수 \checkmark / \texttimes
\end{itemize}

\textbf{시민 참여}
\begin{itemize}
    \item 공청회/의견 조회 실시 여부 \checkmark / \texttimes
    \item 시민 감사단 구성 및 운영 여부 \checkmark / \texttimes
    \item 알고리즘 자문위원회 운영 여부 \checkmark / \texttimes
    \item 시민 만족도 조사 정기 실시 \checkmark / \texttimes
\end{itemize}

\textbf{기술적 건전성}
\begin{itemize}
    \item 성능 모니터링 체계 구축 \checkmark / \texttimes
    \item 드리프트 탐지 메커니즘 존재 \checkmark / \texttimes
    \item 비상 시 수동 전환 절차 수립 \checkmark / \texttimes
    \item 데이터 계보(Lineage) 추적 가능 \checkmark / \texttimes
\end{itemize}

\textbf{품질 인증}
\begin{itemize}
    \item NIA AI 품질 인증 취득 여부 \checkmark / \texttimes
    \item TTA AI 신뢰성 인증 취득 여부 \checkmark / \texttimes
    \item 정기 품질 재평가 일정 수립 \checkmark / \texttimes
\end{itemize}
\end{practicaltool}

%----------------------------------------------------------
\subsection*{제18장 요약}

본 장에서는 공공부문 AI 거버넌스의 특수성과 실무 체계를 다루었다. 공공 AI는 단순 효율성을 넘어 형평성과 민주적 책임성이라는 가치를 구현해야 하며, 네덜란드 SyRI와 호주 Robodebt 사례는 이러한 원칙을 위반했을 때의 심각한 결과를 보여준다. 행정기본법 제20조의 자동적 처분 유형별 자동화 허용 범위를 분석하고, 공공데이터의 AI 활용 시 가명처리, 데이터 결합, 마이데이터 등의 거버넌스 요건을 검토하였다. 알고리즘 투명성 레지스터(네덜란드/영국 사례)와 시민 참여 메커니즘(공청회, 시민 감사단, 알고리즘 자문위원회)을 통한 민주적 통제 체계를 제시하였으며, 공공 AI 조달 시 RFP 작성, 평가 기준, 계약 조건의 실무 지침을 상세히 다루었다. NIA/TTA 품질 인증 체계와 공공 AI 영향평가 자동화 도구를 소개하고, 교육 AI와 민원상담 챗봇의 거버넌스 사례를 통해 공공 AI가 시민의 권리를 보호하면서도 행정 효율성을 달성할 수 있음을 확인하였다.
